\songtitle{Trois pavillons sur le quai}

\setin{1761/1797}
\tag*{2003}\tag{folk-ballad}\tag{cajun}\tag{french-lyrics}\tag{original-source}\tag{british-louisiana}
\begin{notes}
Lyrics by ChatGPT based on scenes of \emph{British Louisiana} (2003), an alternate history by Carlos Thompson (Part 2, \emph{Loyalist Louisiana}, and Part 10, \emph{The 1795 War}) -- a salt merchant on the New Orleans quay from 1761 to 1797, under the French, British and Spanish flags, and the British guns. Retrieved from Wyomee Dank's Suno page (V6-MINI, 24 November 2026).
\end{notes}
\begin{style}
Cajun ballad, close solo lead voice in a high, light upper register with clipped, rhythmic staccato phrasing and Southern Louisiana Cajun French diction; an unaccompanied voice carries a narrow, lilting melody over one open-string drone, while one fiddle plays short bowed answers between stanzas; tape warmth and narrow frequency range; gentle two-beat sway, with the same verse tune lifting slightly in each refrain and the final refrain left unaccompanied.
\end{style}
\begin{lyrics}
\songpart{Verse 1 - Solo}
Du temps du roi de France,\\
Mon père vendait du sel au quai.\\
Il m’a laissé son gros registre\\
Pis les noms de ceux qui lui devaient.

\songpart{Verse 2 - Solo}
Quand les Anglais ont pris la ville,\\
Leur commis m’a fait signe d’entrer.\\
Il a écrit mon nom à sa manière.\\
Moi, j’ai signé comme on m’a montré.

\annotation{Refrain - Solo}\\
Trois pavillons au bout du quai,\\
Pis j’ai la même clef dans ma main.\\
Quel roi c’est qui va se lever\\
Pour ouvrir ma porte demain ?

\annotation{Short fiddle interlude}

\songpart{Verse 3 - Solo}
Quand les Espagnols ont pris la ville,\\
J’ai vu leurs soldats sur le quai.\\
Y en a un qui m’a demandé du sel.\\
J’ai dit qu’il paierait le mois d’après.

\annotation{Refrain - Solo}\\
Trois pavillons au bout du quai,\\
Pis j’ai la même clef dans ma main.\\
Quel roi c’est qui va se lever\\
Pour ouvrir ma porte demain ?

\annotation{Short fiddle interlude}

\songpart{Verse 4 - Solo}
Asteur les Anglais nous assiègent.\\
Du fond d’la boutique, j’entends l’canon.\\
Ma fille me demande qui commande icitte.\\
J’lui dis : « Viens donc écrire ton nom. »

\annotation{Final Refrain - Unaccompanied solo}\\
Trois pavillons au bout du quai,\\
Pis j’ai la même clef dans ma main.\\
Quel roi c’est qui va se lever\\
Pour ouvrir ma porte demain ?
\end{lyrics}

